\section{Conclusiones}

Las conclusiones de esta ICR deberán responder directamente a las preguntas planteadas en la introducción. Para mantener coherencia con el alcance del estudio, el cierre no buscará generalizar sobre toda la vivienda hñähñu del Valle del Mezquital, sino explicar qué revelan cuatro casos seleccionados dentro de un registro más amplio sobre sistemas constructivos, materiales locales, transformaciones y saberes asociados.

\subsection{Respuesta a las preguntas de investigación}

El cierre se organizará en torno a cinco respuestas:

\begin{enumerate}
	\item qué sistemas constructivos fueron identificados y cómo pueden clasificarse;
	\item qué relación existe entre materiales, disponibilidad local y decisiones constructivas;
	\item qué aportó la cartografía para comprender la ubicación de las construcciones en el territorio;
	\item qué cambios y permanencias se reconocieron mediante fotografías y relatos;
	\item qué estructura de ficha, síntesis comparativa y criterios de documentación permite ordenar el diagnóstico.
\end{enumerate}

Esta organización permitirá que las conclusiones se desprendan del diagnóstico y no repitan el marco teórico.

\subsection{Aportes de la investigación}

Los aportes esperados son cuatro. Primero, un registro descriptivo de construcciones tradicionales poco documentadas. Segundo, una tipología inicial de sistemas constructivos persistentes en dos comunidades. Tercero, una base cartográfica y testimonial que permita discutir su valor como patrimonio vernáculo biocultural. Cuarto, un conjunto de fichas y criterios básicos para orientar la documentación posterior.

La contribución principal no estará en proponer una intervención arquitectónica, sino en demostrar que estas construcciones todavía contienen información técnica, material y cultural suficiente para justificar su documentación sistemática.

\subsection{Límites y líneas de continuidad}

El estudio reconoce límites claros: número reducido de casos, ausencia de levantamientos arquitectónicos medidos, imposibilidad de caracterizar en laboratorio todos los materiales y necesidad de proteger información sensible. Estos límites no invalidan el diagnóstico; delimitan su alcance.

Como líneas de continuidad se plantean la ampliación del registro a más comunidades, la caracterización material de la piedra local denominada ``tepetate'', la comparación temporal con nuevos registros fotográficos y la incorporación progresiva de nuevos sistemas constructivos documentados con la misma ficha base.
